% !TeX program = pdflatex
% ─────────────────────────────────────────────────────────────────────────────
%  papers/estatistica.tex — Estatística: o espaço de probabilidade e a
%  realização da onda.
%
%  Compila:  cd papers && pdflatex estatistica.tex
% ─────────────────────────────────────────────────────────────────────────────
\documentclass[11pt,a4paper]{article}
\usepackage[utf8]{inputenc}\usepackage[T1]{fontenc}\usepackage[brazilian]{babel}
\usepackage{amsmath,amssymb,amsthm}
\usepackage[margin=2.4cm]{geometry}
\usepackage{booktabs}\usepackage{array}\usepackage{enumitem}\usepackage{xcolor}
\usepackage[htt]{hyphenat}
\usepackage[hidelinks,breaklinks]{hyperref}
\usepackage{microtype}

\sloppy
\emergencystretch=2em

\definecolor{cinza}{gray}{0.35}
\newcommand{\code}[1]{\texttt{\small #1}}
\newcommand{\caixa}[1]{%
  \par\smallskip\noindent
  \begin{center}
  \fbox{\begin{minipage}{0.94\textwidth}\centering\smallskip #1\smallskip\end{minipage}}
  \end{center}\smallskip}
\newcommand{\abs}[1]{\lvert #1\rvert}
\newcommand{\medido}[1]{\par\medskip\noindent{\small\color{cinza}\textbf{Medido.} #1}\par\medskip}
\providecommand{\interfacen}[1]{}
\newtheorem{teorema}{Teorema}
\newtheorem{lema}[teorema]{Lema}
\newtheorem{definicao}[teorema]{Definição}
\newtheorem{corolario}[teorema]{Corolário}
\newtheorem{proposicao}[teorema]{Proposição}
\newtheorem{observacao}[teorema]{Observação}
\newtheorem{problema}[teorema]{Problema}
\renewcommand{\proofname}{Demonstração}

% ─────────────────────────────────────────────────────────────────────────────
%  papers/gkcapa.tex — a capa do Reino, mínima, para os papers em `article`.
%
%  O estilo.tex completo é do `report` (títulos de capítulo, skak, …). Aqui fica
%  só o que a capa precisa, para o paper continuar a compilar sozinho com
%  pdflatex. O tradutor próprio (tests/tex) carrega o estilo.tex da raiz / o
%  symlink papers/estilo.tex e avalia o mesmo `\gkcapa`.
% ─────────────────────────────────────────────────────────────────────────────
\definecolor{ouro}{HTML}{B8912F}
\definecolor{ouroclaro}{HTML}{F3E9CE}
\definecolor{tinta}{HTML}{1E1B16}
\definecolor{regua}{HTML}{6E6858}
\providecommand{\gknota}{\fontsize{7.62}{11.05}\selectfont}
\providecommand{\gktexto}{\fontsize{10.50}{15.22}\selectfont}
\providecommand{\gksub}{\fontsize{12.33}{17.87}\selectfont}
\providecommand{\gksec}{\fontsize{14.47}{20.98}\selectfont}
\providecommand{\gkcap}{\fontsize{16.99}{24.63}\selectfont}
\providecommand{\gktit}{\fontsize{23.42}{33.95}\selectfont}
\providecommand{\Prj}{\mathbb{P}}
\providecommand{\Trf}{\mathcal{T}}
\providecommand{\gkcapa}[3]{%
\title{\vspace{-0.6cm}
{\color{ouro}\rule{\textwidth}{1.4pt}}\\[6mm]
\textbf{\gktit\color{tinta} Reino Dourado}\\[3mm]{\gksec\itshape\color{regua} Golden Kingdom}\\[8mm]
{\gkcap\scshape\color{ouro} #1}\\[5mm]
{\gksub\color{regua} #2}\\[2mm]
{\gktexto\color{regua}\itshape #3}\\[7mm]
{\color{ouro}\rule{\textwidth}{0.6pt}}\\[9mm]{\gknota\color{regua}\begin{minipage}{\textwidth}\setlength{\parindent}{0pt}%
\textbf{Aviso de Isenção Algorítmica:} O universo, as leis e as entidades contidas nestas páginas ---
incluindo felinos matriciais, aves de rapina algébricas e amuletos de controle de fluxo --- são
construções de pura ficção formal. Esta obra não descreve a realidade física, não serve como
engenharia reversa do cosmos e não constitui doutrina religiosa. Qualquer semelhança com bugs reais do seu
sistema operacional é mera coincidência (ou falta de reza). Prossiga por sua própria conta e
risco.\end{minipage}}}
\author{}\date{}}


\begin{document}
\interfacen{16}
\gkcapa{Estatística --- a probabilidade como interface da realização}
       {multiplicidade $\to$ módulo; dualidade $\to$ fase}
       {a matéria é uma onda de probabilidades}
\maketitle

\begin{center}
\small
\textit{O que se afirma: esta obra constrói os seus objetos e mede-os.
Cada número publicado tem o programa que o produziu, e cada enunciado tem a sua demonstração
--- a bateria corre com um comando e diz o que mediu.
Quem quiser conferir não precisa de acreditar em nada do que aqui está escrito.}
\end{center}

\vspace{1.2em}

\section*{A tese desta parte}

\begin{center}
\large
\textbf{a matéria é uma onda de probabilidades}
\end{center}

\vspace{0.4em}
\begin{center}
\small
(esta é a leitura final da realização; \emph{não} é a definição inicial de $\psi$)
\end{center}

\vspace{0.6em}

A regra de Born aparece na Quântica, mas não está localizada no Corpo de Partículas
(Obs.~0.215). A probabilidade-contagem $p_G=G/|I|$ já está no palco desde a realização
(Def.~0.23; Obs.~0.214). O espaço de probabilidade não é apenas o lugar onde se contam
resultados: é o suporte onde a matéria se realiza estatisticamente.

A pergunta correcta deixa de ser «como aplicar Born?» e passa a ser:

\begin{center}
\textit{como a função de onda nasce da estrutura que já temos?}
\end{center}

A onda não se constrói a partir de $\sqrt{G/|I|}$. Constrói-se de baixo para cima:

\begin{center}
Lei 0/7 $\to$ cisão $\to$ duas leituras $\to$ involução no par $(x,y)$\\[2pt]
$d(xy)$ $\to$ $\int f+\int f^{-1}$\\[2pt]
cisão por paridade $\to$ $\exp(tJ)=c+sJ$ $\to$ $c^2+s^2=1$ $\to$ $\theta$, $\pi$\\[2pt]
$x^2=x+1$ $\to$ $\varphi$ $\to$ $A_n$\\[2pt]
$\psi_n=A_n\,e^{J\theta_n}$
\end{center}

Neste quadro, a probabilidade não entra como princípio primeiro. O que se constrói primeiro é a realização; depois a representação adequada; depois o peso da fibra; depois a normalização; depois a distribuição de contagem; e só então a leitura observável que se exprime como $|\psi|^2$. A ordem correta é
\[
\boxed{
\text{estrutura}
\longrightarrow
\text{representação}
\longrightarrow
\text{peso}
\longrightarrow
\text{normalização}
\longrightarrow
\text{distribuição}
\longrightarrow
\text{Born}
}
\]
Aí está a distinção que o presente texto precisa manter: a distribuição não é a origem da onda; ela é o resultado da normalização de um peso já construído pela estrutura. A métrica não produz a distribuição. Ela fornece uma leitura geométrica da estrutura de concordância do suporte. A multiplicidade $G$, obtida pela contagem das dobras, fornece o peso; a sua normalização $G/|I|$ produz a distribuição; e Born é a leitura física dessa distribuição.

Em termos precisos,
\[
\boxed{\text{métrica} \neq \text{probabilidade}}
\]
\[
\boxed{\text{distribuição} \neq \text{regra de Born}}.
\]
A primeira é o objeto matemático que emerge da contagem e da normalização; a segunda é a leitura física que se faz sobre esse objeto. O papel de Born não é fabricar a distribuição; é interpretá-la. Em outras palavras, a estrutura não é produzida por Born, e Born não é a primeira camada da teoria. Ele é a última etapa do enquadramento estatístico, depois que a geometria, a fase e a multiplicidade já ficaram em lugar.

A probabilidade é leitura estatística dessa realização, não a receita.
Em termos rigorosos, a ordem é
\[
\text{realização} \;\to\; \psi=Ae^{J\theta} \;\to\; |\psi|^2\to P_U.
\]
Antes da ponte de Born, $\psi$ é uma \emph{realização de onda candidata}: geométrica e dinâmica,
mas ainda não uma distribuição de probabilidade. A frase “a matéria é uma onda de probabilidades”
refere-se à interpretação final do mesmo objecto, não ao ponto de partida da construção.

\vspace{1em}
\begin{center}
\renewcommand{\arraystretch}{1.25}
\begin{tabular}{@{}p{3.4cm}p{7.6cm}p{3.6cm}@{}}
\toprule
\textbf{no espaço} & \textbf{é, nesta construção} & \textbf{onde} \\
\midrule
ocorrências / dobra &
lista $I$, realização $\pi$, multiplicidade $G$ &
Def.~0.23 \\
medida de contagem &
$p=G/|I|$ (estatística, não fonte do módulo) &
Obs.~0.214; Teor.~0.44 \\
projeção &
$x^2=x$ (idempotente / estado puro) &
Obs.~0.215 \\
fase / rotação &
$x^2=-1$ (rotor); degraus; $\operatorname{Res}_\theta=0$ &
Teor.~0.102(2); 0.22 \\
crescimento / amplitude &
$x^2=x+1$ (estrela / espiral / Alonzo) &
Teor.~0.102; catálogo 0.20.2 \\
costura norma &
exterior + rotor: $a^2\to a^2+b^2$ &
catálogo (ext.+rotor) \\
função de onda &
objecto em que as três equações se realizam juntas &
Def.~\ref{def:psi} \\
\bottomrule
\end{tabular}
\end{center}

% ============================================================
\section{Porque esta parte vem aqui}
% ============================================================

A Parte Quântica (XII) organizou o modo espectral --- massa $|\log\rho|$, luz $\rho=1$ ---
sem magma nova. Lá se leu explicitamente:

\begin{itemize}[leftmargin=1.4em,itemsep=2pt]
  \item $P=k/N$ já realizado (contagem);
  \item $p_i=|\langle i|\psi\rangle|^2$ continua \emph{não localizada} (Obs.~0.215);
  \item Parseval da dourada conserva a norma $L^2$ da transformada, \emph{não} produz Born
        (Obs.~0.220);
  \item GLH, dobras e Lei~7 realizam contagem e conversão no palco, \emph{não} Born
        (Obs.~0.221).
\end{itemize}

Born só sai com estatística. Mas estatística, nesta construção, não começa por um axioma
de medida abstracta: começa pela multiplicidade que a realização já conta.

O Resumo da Física já declara o eixo dual que esta Parte vai usar:

\begin{quote}
\small
«o cartesiano e o polar, onde as fases somam e os módulos multiplicam.
E as duas fecham: batendo alternadas, os intervalos encaixam e o processo termina
numa igualdade --- o ponto onde para é o ponto fixo que as duas faces passam a partilhar.»
\end{quote}

A transformada já construída é pura fase (Teor.~0.160(1)): os caracteres $\chi_k:X\to\{\pm1\}$
não carregam magnitude. Toda a magnitude do transformado vem do campo $G$, e $G$ conta
dobras --- mas a contagem regista a dobra, não fabrica sozinha a amplitude da espiral.
O que falta para a onda é a realização conjunta das três equações
(projeção, rotação, crescimento) no mesmo objecto, com teste de resíduo no lado rotacional
e fecho ainda aberto no lado hiperbólico.

% ============================================================
\section{Da multiplicidade à medida}
% ============================================================

\begin{definicao}[realização e multiplicidade --- já construídas]
\label{def:realizacao}
Seja $I$ a lista finita (Teor.~0.20) e $X$ um conjunto com igualdade decidível.
Uma realização é $\pi:I\to X$. A multiplicidade geométrica é
\[
G(x)=\bigl|\pi^{-1}(x)\bigr|=\bigl|\{i\in I:\pi(i)=x\}\bigr|.
\]
O suporte de $G$ é $\operatorname{supp}G=\operatorname{im}\pi$ (Def.~0.23).
\end{definicao}

\begin{definicao}[espaço de probabilidade clássico]
\label{def:espaco-prob}
O espaço de probabilidade clássico associado à realização é o triplo
\[
(\Omega,\mathcal{F},P)
\quad\text{com}\quad
\Omega=\operatorname{supp}G,\qquad
\mathcal{F}=2^\Omega,\qquad
P(A)=\frac{1}{|I|}\sum_{x\in A}G(x)=\frac{|A|_G}{|I|},
\]
onde $|A|_G=\sum_{x\in A}G(x)$ é a contagem com multiplicidade.
A ordem estrutural é a seguinte: primeiro há a multiplicidade $G$, que é o peso da fibra; depois a normalização pela massa total $|I|$; depois a distribuição $p(x)=G(x)/|I|$; e só então a leitura estatística $P$. Não se introduz medida nova nem regra de Born para criar a distribuição: a distribuição emerge da estrutura e a regra de Born passa a ser a sua interpretação física. Quando se identifica a fibra com a célula, $P(\{x\})=G(x)/|I|$.
É a mesma contagem de Hurwitz (Teor.~0.44) lida como probabilidade.
Não se introduz medida nova. A finitude fecha.
\end{definicao}

% ============================================================
\section{A ponte: projeção, rotação e crescimento da espiral}
\label{sec:ponte}
% ============================================================

A distribuição de contagem $p=G/|I|$ regista a dobra da realização.
Heighway e a espiral quadrada podem partilhar o mesmo $|I|$ e ter multiplicidades
radicalmente diferentes: a medida sozinha \emph{não} distingue realizações e
\emph{não} fabrica, por si, a amplitude geométrica.
Não se fecha portanto
\[
\psi=\sqrt{G/|I|}\,e^{i\theta}
\]
como se as duas fontes fossem independentes e fundamentais.
Inverte-se a pergunta.

\begin{problema}
\label{prob:central}
Dada a espiral $(z_n)$ construída no Reino, qual é a realização de onda?
É $\psi_n\equiv z_n/\|z_n\|$?
Ou $\psi_n=A_n\,e^{i\theta_n}$ com $A_n$ produzido pela própria recorrência?
Nesse regime, $G/|I|$ pode aparecer \emph{depois}, como estatística observável
da realização --- não como fonte primária da amplitude.
\end{problema}

\subsection{Três equações, três papéis}

Os documentos separam com clareza três equações e não as fundem:

\begin{center}
\renewcommand{\arraystretch}{1.25}
\begin{tabular}{@{}llll@{}}
\toprule
equação & tipo & papel & sítio \\
\midrule
$x^2=x$ & idempotente & projeção / estado puro & Obs.~0.215; projector $\rho$ \\
$x^2=-1$ & rotor (elíptico, ord.~4) & fase / rotação & Teor.~0.102(2); degrau $w=2$ \\
$x^2=x+1$ & estrela (hiperbólica) & crescimento da espiral & Teor.~0.102; Alonzo; $\varphi$ \\
\bottomrule
\end{tabular}
\end{center}

\begin{itemize}[leftmargin=1.3em,itemsep=2pt]
  \item $x^2=x$ é o projector. Não é rotação nem expansão. É a realização de um estado puro.
  \item $x^2=-1$ é o rotor: ordem finita 4, elíptico, ciclo
        $1\to x\to-1\to-x\to1$. É a fase / orientação no plano.
  \item $x^2=x+1$ é a estrela (Star(K)): hiperbólica, sem ordem finita; as raízes são
        $\varphi$ e $\psi$. É o crescimento / deformação da espiral.
\end{itemize}

A onda aparece quando estas três leituras se realizam no \emph{mesmo objecto},
não quando se anexa uma amplitude estatística a uma fase.

\subsection{Exterior, rotor e a costura que produz a norma}

O catálogo regista a costura explícita:

\begin{itemize}[leftmargin=1.3em,itemsep=2pt]
  \item \textbf{Exterior sozinho}: $a+be$, $e^2=0$, norma $a^2$ (cega a $e$),
        assinatura com divisor de zero --- o cone nulo.
  \item \textbf{Exterior + rotor}: troca $e^2=0$ por $i^2=-1$. Resultado $\mathbb{C}$,
        norma $N=a^2+b^2$, equação $x^2=-1$.
\end{itemize}

O exterior sozinho é cone; com o rotor vira corpo. A norma passa de $a^2$ para $a^2+b^2$.
Isso é a costura que o documento já fechou (forma $(2,0,0)$, ISO-DUAL).
A fase (rotação) e a estrutura que sustenta o módulo (norma definida) nascem juntos
nessa substituição --- não de uma contagem posterior.

\subsection{A espiral de ouro e o cone nulo}

No catálogo (0.20.2) a espiral áurea é o crescimento sequencial que selecciona $\varphi$
(ângulo $2\pi/\varphi^2$). Uma dimensão acima, a malha erguida é um \textbf{cone nulo}:
a espiral plana é a sombra da estrutura que vive no cone.
O dual da espiral áurea é o cone nulo; o teste \texttt{tests/cone\_nulo.c} regista-o.

Alonzo (corpo canónico da Parte Fractal) realiza a composição; $\varphi$ e $\psi$ são
as raízes de $x^2=x+1$; o dual Cantor~$\leftrightarrow$~Julia liga os dois tecidos.
A forma áurea $\varphi\cdot\varphi=\varphi+1$ é o produto que é a soma --- crescimento
escrito como equação de ponto fixo.

\subsection{Degraus da fase: o lado rotacional ainda vale}

Os degraus (Teor.~0.22) continuam a dar a torre rotacional discreta:
\[
\Delta\theta_w=\frac{2\pi}{2^w},\qquad
\operatorname{ord}(\rho_w)=2^w,\qquad
\sum_{k=0}^{2^w-1}\Delta\theta_w=2\pi.
\]
No caso discreto, o resíduo de ciclo é
\[
\operatorname{Res}^{\mathrm{ciclo}}_\theta
=\sum_{k=0}^{2^w-1}\Delta\theta_w-2\pi=0.
\]
Esse é o teste de fecho da \emph{parte fase} no rotor. Os degraus não fabricam o módulo;
 fabricam a orientação. O módulo, se vier da espiral, vem da recorrência hiperbólica
($x^2=x+1$), não de $\sqrt{G/|I|}$.

Na espiral áurea a etapa angular é
\[
\Delta\theta=\frac{2\pi}{\varphi^2},
\]
mas como $\varphi^2$ é irracional, não existe $N\in\mathbb N$ tal que $N\Delta\theta=2\pi$.
Logo, para a espiral, o fechamento plano não é periódico finito. O objecto relevante é a
\emph{identidade local do passo angular}
\[
\operatorname{Res}^{\mathrm{passo}}_\theta
=\Delta\theta-\frac{2\pi}{\varphi^2}=0,
\]
não um fechamento de ciclo inteiro no plano. Isso separa claramente:
\[
\text{rotor discreto} = \text{fecho periódico};\qquad
\text{espiral áurea} = \text{passo angular exacto e não periódico.}
\]

\begin{teorema}[fecho rotacional --- resíduo nulo]
\label{thm:fecho-fase}
Sobre o ciclo gerado por $\rho_w$, $\sum\Delta\theta_w=2\pi$ e
$\operatorname{Res}^{\mathrm{ciclo}}_\theta=0$.
A condição está embutida na ordem do gerador (Teor.~0.22). Isso valida a parte
rotacional da ponte; não valida ainda a origem da amplitude nem o fechamento plano da espiral áurea.
\end{teorema}

\subsection{Forma diferencial do Teorema Central}

O Teorema Central não começa pela integral: começa pelo diferencial do produto.
Seja
\[
y=f(x),\qquad x=f^{-1}(y).
\]
O termo central da identidade é
\[
\boxed{y\,dx+x\,dy=d(xy)}.
\]
Substituindo $y=f(x)$:
\[
\boxed{f(x)\,dx+f^{-1}(y)\,dy=d\bigl(x\,f(x)\bigr)},
\]
com $dy=f'(x)\,dx$. Esta é a forma diferencial do Teorema Central.

A forma integral é a leitura acumulada: integrando de $a$ a $b$ no domínio e de $f(a)$ a $f(b)$ na imagem,
\[
\int_a^b y\,dx+\int_{f(a)}^{f(b)}x\,dy=\bigl[xy\bigr]_a^b,
\]
isto é
\[
\boxed{\int_a^b f(x)\,dx+\int_{f(a)}^{f(b)}f^{-1}(y)\,dy=b\,f(b)-a\,f(a).}
\]
É a partição reversível do mesmo rectângulo pelos dois cortes --- domínio (Riemann/Hurwitz)
e imagem (Lebesgue/Gentil) --- sem resto. O catálogo já descreve a soma das duas leituras
como essa partição; a forma diferencial mostra que a identidade é originalmente
o diferencial do produto, e a integral é a sua acumulação.

\subsection{Lei 0, Lei 7, Lei 2 --- e o que cada uma entrega}

Antes de especializar $f'=f^{-1}$, a derivação tem de ser rastreada:

\begin{center}
\renewcommand{\arraystretch}{1.25}
\begin{tabular}{@{}llp{7.5cm}@{}}
\toprule
lei & conteúdo nos docs & o que entrega à ponte \\
\midrule
Lei 0 & $(+1)\oplus(-1)$; $0^\dagger=\infty$ &
operação fundamental nos extremos \\
Lei 7 & $H\times H^*$; dois tecidos ligam sem fundir &
transporte / ligação sem fusão \\
vinco & $1=\mathrm{cl}(L_0,L_7)$ &
encontro das duas \\
retorno & $\pi_U(F(L_7))=L_0$, $\mathrm{Res}=0$ &
ciclo fechado (Teor.~univ:55) \\
Lei 2 & $K^{**}=K$; catálogo: «$f\leftrightarrow f^{-1}$ (a Lei 2)» &
involução dual explicitamente identificada \\
\bottomrule
\end{tabular}
\end{center}

Lei 0 + Lei 7 fecham o ciclo (vinco + retorno). A forma diferencial do produto
\[
y\,dx+x\,dy=d(xy)
\]
e a sua acumulação integral são o Teorema Central --- partição reversível do rectângulo
pelos dois cortes. Isso não é lei nova.

A identificação $f\leftrightarrow f^{-1}$ está no catálogo como conteúdo da \textbf{Lei 2},
não como teorema derivado passo a passo de Lei 0+7. Enquanto essa derivação explícita
não estiver escrita com resíduo zero, \emph{não} se declara $f'=f^{-1}$ ``obrigatória só porque
$K^{**}=K$''. Declara-se: a involução dual está nomeada na Lei 2; o pedido analítico
$f'=f^{-1}$ é a realização dessa involução \emph{dentro da família de potências}.

\paragraph{Justificativa da casa: sem Lei 8 e com retorno de 7 a 0.}
O repositório toma como afirmação canónica que as leis fecham em $0\ldots 7$, e que não existe
uma \textbf{Lei 8} como nova lei dinâmica. Em termos do documento de referência do sistema,
$8$ é a cardinalidade do conjunto das leis, não uma nova etapa da dinâmica; a Lei 7 não é uma
parede, mas um corpo novo (o octonião dual) que reaplica a mesma estrutura e reentra no ciclo.
Este ponto é explícito no \texttt{memoria/project-checkpoint-2026-08-09-oito-leis.md} e no \texttt{papers/aranha.tex}: o ciclo da aranha escreve, lê, decide e fecha, sem introduzir um
sétimo passo extra. O fechamento é a volta ao ponto de origem, isto é, a descida $7\to 0$ pelo
próprio ciclo, e não a criação de uma nova lei. A regra da casa é, portanto: 
\[
\boxed{\text{Lei 7 reaplica a estrutura e fecha em Lei 0, sem Lei 8}.}
\]
O mesmo padrão aparece no ciclo da aranha em \texttt{papers/aranha.tex}: escrita, leitura,
decisão e fecho formam um retorno estrutural ao ponto base e não uma lei suplementar.

\subsection{Unicidade do expoente --- não unicidade global da potência}

Na família monomial $f(x)=a\,x^b$, a inversa é
\[
 f^{-1}(x)=\left(\frac{x}{a}\right)^{1/b}
\]
quando $a>0$ e $b\neq 0$. A derivada é
\[
 f'(x)=ab\,x^{b-1}.
\]
A condição $f'(x)=f^{-1}(x)$ obriga, em primeiro lugar, à igualdade dos expoentes:
\[
 b-1=\frac1b
\qquad\Longrightarrow\qquad
b^2-b-1=0
\qquad\Longrightarrow\qquad
b=\varphi.
\]
Mas a igualdade dos coeficientes impõe também
\[
 ab=a^{-1/b}
\qquad\Longrightarrow\qquad
a^{1+1/b}=\frac1b.
\]
Para $b=\varphi$, isso fixa o coeficiente
\[
 a=\varphi^{-1/\varphi}.
\]
Assim, dentro da família $a x^b$, a condição $f'=f^{-1}$ força \textbf{não só o expoente} $b=\varphi$, mas também o coeficiente $a$.

Controlo negativo (\texttt{tests/aurea.c}): se $b\neq\varphi$, a igualdade quebra e o desvio
acompanha $|b^2-b-1|$. Unicidade do expoente \emph{dentro} das potências: sim, com resíduo 0.

O catálogo \textbf{não} afirma que as potências sejam as únicas funções $C^1$ que satisfazem
$f'=f^{-1}$ no contínuo. Não se promove essa unicidade global. O quadro de trabalho é
a família potência; dentro dela o ouro é forçado, e o coeficiente também fica fixado.
\[
f(x)\,dx+x\,f^{-1}(x)\,dx=d\bigl(x\,f(x)\bigr)
\]
(o factor $x$ no segundo termo é essencial).

A matriz áurea
\[
M_\varphi=\begin{pmatrix}1&1\\1&0\end{pmatrix}
\]
tem $\det=-1$, autovalores $\varphi$ e $-1/\varphi$: uma direcção expande, a conjugada
contrai e inverte a orientação. O sinal negativo da folha conjugada é esse determinante,
não um postulado extra.

Cadeia (com o que está fechado e o que ainda é identificação de catálogo):
\begin{center}
\renewcommand{\arraystretch}{1.3}
\begin{tabular}{@{}cl@{}}
Lei 0 + Lei 7 & vinco + retorno ($\mathrm{Res}=0$) \\
$y\,dx+x\,dy=d(xy)$ & forma diferencial do Teorema Central \\
$\displaystyle\int f+\int f^{-1}=bf(b)-af(a)$ & forma integral (acumulação) \\
Lei 2 (catálogo) & nomeia $f\leftrightarrow f^{-1}$ \\
$f'=f^{-1}$ em potências & realização analítica no quadro $ax^b$ \\
$b=\varphi$ & expoente forçado (unicidade \emph{dentro} das potências) \\
\end{tabular}
\end{center}

\subsection{Gentil--Hurwitz--Lebesgue: os dois lados da mesma medida}

O teorema central (Teor.~cat:0.7 / thm:central) lê os dois lados pela medida:

\begin{itemize}[leftmargin=1.2em,itemsep=1pt]
  \item \textbf{Hurwitz} conta (discreto, esfera, bilinear $1,2,4,8$);
  \item \textbf{Lebesgue} transporta a ordem ($\mathbb{Q}\to\mathbb{R}$);
  \item \textbf{Gentil} integra (contínuo, hipérbole, sem limite).
\end{itemize}

A bijeção dual é a estrela, realizada pela forma integral acima.
Involução $f\leftrightarrow f^{-1}$ (Lei~2). Nunca se avalia uma raiz: os dois cortes
dão a mesma contagem.

Na leitura logarítmica a função áurea lineariza:
\[
\ln f(e^u)=\varphi\,u+\ln a,
\]
e o período espectral das voltas da reescala é
\[
\tau_0=\frac{2\pi}{\ln\varphi}.
\]
Recorrência $\varphi^n$ produz amplitude; $\ln\varphi$ produz a contagem das voltas (fase).
São duas projecções da mesma máquina --- a que satisfaz $f'=f^{-1}$ --- não dois
ingredientes colados.

\subsection{Dois cortes, a mesma conservação}

A arquitectura comum aos dois sítios:

\begin{center}
objeto $=$ face$_1$ $+$ face$_2$
\qquad;\qquad
reflexão/troca $\Longrightarrow$ conservação
\end{center}

\textbf{Corte 1 --- Teorema Central (Teor.~0.44).}
Domínio $\mid$ imagem. A involução actua no \emph{par}:
\[
(x,y)\;\overset{\nu}{\longmapsto}\;(y,x),\qquad \nu^2=\mathrm{id}.
\]
Forma infinitesimal:
\[
\underbrace{y\,dx}_{\text{domínio}}+\underbrace{x\,dy}_{\text{imagem}}=\underbrace{d(xy)}_{\text{vinco/fecho}}.
\]
Acumulação:
\[
\int_a^b f+\int_{f(a)}^{f(b)}f^{-1}=bf(b)-af(a).
\]
A notação $f\leftrightarrow f^{-1}$ é coordenatização do par quando se escolhe a carta $y=f(x)$
--- não um mapa novo $V\to V$ (catálogo e Teor.~0.44(3)).

\textbf{Corte 2 --- exponencial / forma polar.}
Par $\mid$ ímpar. Com $J^2=-1$ (segundo degrau da largura, ordem 4):
\[
\exp(tJ)=c(t)+s(t)J,
\]
\[
c(t)=\sum_k\frac{(-1)^k t^{2k}}{(2k)!},\qquad
s(t)=\sum_k\frac{(-1)^k t^{2k+1}}{(2k+1)!}.
\]
Conservação pela reflexão:
\[
(c+sJ)(c-sJ)=c^2+s^2=1.
\]
A exponencial é o dicionário entre as duas faces (módulo multiplica, fase soma);
a forma polar não é importada --- é o Teor.~0.10(4)/0.22 onde as faces convivem.
Julia ($z\mapsto z^2$) e Cantor ($\theta\mapsto 2\theta$) conjugam-se por $z=e^{2\pi i\theta}$.

\textbf{$\pi$ produzido pelo fluxo, não postulada.}
Uma vez fixada a parametrização do fluxo (por radianos na carta do rotor), escreve-se
\[
\pi=\min\{t>0:\exp(tJ)=-1\}.
\]
Quarto de volta $\to J$, meia volta $\to-1$, volta inteira $\to 1$.
A fase é coordenada do fluxo; não se inventa a partir de $\pm1$.

\subsection{A realização de onda --- de baixo para cima}

\begin{definicao}[realização de onda]
\label{def:psi}
A realização de onda é o objecto em que as cisões e as recorrências já construídas se lêem juntas:
\begin{enumerate}[leftmargin=1.2em,itemsep=2pt]
  \item \textbf{Cisão / duas leituras} (Lei 0, Lei 7, Teor.~0.44): par $(x,y)$, involução $\nu$,
        forma $d(xy)$.
  \item \textbf{Rotação / fase} ($J^2=-1$, degraus, $\exp(tJ)=c+sJ$):
        $\theta$ coordenada do fluxo, com $\operatorname{Res}^{\mathrm{ciclo}}_\theta=0$ no rotor discreto
        e $\operatorname{Res}^{\mathrm{passo}}_\theta=0$ na espiral áurea.
  \item \textbf{Crescimento / amplitude} ($x^2=x+1$, $f'=f^{-1}$ em potências):
        $A_{n+1}=\varphi A_n$, forçado no quadro $ax^b$ (catálogo §0.64).
\end{enumerate}
Escreve-se
\[
\boxed{\psi_n=A_n\,e^{J\theta_n}}
\]
com $A_n$ da recorrência e $e^{J\theta_n}$ da cisão/rotação --- duas leituras da mesma realização,
não dois ingredientes colados. No fluxo puro $\exp(tJ)$ tem-se $A\equiv 1$ e $c^2+s^2=1$;
na espiral a recorrência metálica deforma a amplitude. O que se obtém aqui é uma
\emph{realização de onda candidata}: a órbita $T^n\psi_0$ tem forma polar áurea, mas a ponte
estatística para $P_U$ continua aberta e não pode ser usada como hipótese de construção.
Este teorema prova a \emph{consistência} da órbita, não a derivação causal completa de
$T=\varphi e^{J\Delta\theta}$ a partir de $L_0+L_7$. A distinção é importante: o texto
constrói a solução da recorrência e a sua forma polar, mas não declara que essa mesma forma seja
já derivada das leis anteriores; a derivação causal de $T$ permanece como dívida explícita, não como
resultado escondido.
\end{definicao}

\begin{observacao}[probabilidade como interface, não como terceira face]
\label{obs:nao-fecha}
A probabilidade $G/|I|$ é leitura estatística da realização, não a receita que a constrói.
A igualdade $|\psi|^2=P_U$ não é usada para a definir; ela é a ponte de compatibilidade entre as duas leituras do mesmo objecto. Não se promove $\sqrt{G/|I|}$ a amplitude fundamental.
Não se inventa mapa $V\to V$ além do par de leituras (Teor.~0.44(3)).
Em letras garrafais: $\text{norma}\neq\text{probabilidade}\neq\text{amplitude}$.
A estrutura correta é
\[
\text{Geometria}
\;\rightleftarrows\;
P_U
\;\rightleftarrows\;
\text{Dinâmica},
\]
com $P_U$ no centro como interface de leitura, não como terceiro estádio da cadeia. A mesma recorrência que produz crescimento produz a parte radial; o rotor produz a parte angular; a probabilidade não entra para fazer a onda nem para definir a sua forma. A dívida restante é geométrica: mostrar como o par dual
$(\psi^+,\psi^-)$ determina a fibra $\pi_U^{-1}(x)$ e, só então, a contagem $G(x)$.
Fecha: involução no par; $d(xy)$; forma integral; $\exp(tJ)$ com conservação;
$\varphi$ forçado em potências; $\operatorname{Res}^{\mathrm{ciclo}}_\theta=0$ no rotor;
$\pi$ pelo fluxo. O passo que resta operacionalizar sítio a sítio é a identificação explícita
$\psi_n=A_n e^{J\theta_n}$ em Alonzo / espiral de ouro / modo espectral, com o mesmo
critério de resíduo zero, e a posteriori a ponte $P_U\stackrel{?}{=}|\psi|^2$.
\end{observacao}


% ============================================================
\section{Fecho na espiral de ouro}
\label{sec:espiral}
% ============================================================

O catálogo (0.20.2) constrói a espiral de ouro em três passos, sem postular onda:

\begin{enumerate}[leftmargin=1.3em,itemsep=2pt]
  \item crescimento sequencial com máxima simetria (cada semente no maior vão) ---
        selecciona o número mais irracional $\varphi$;
  \item giro entre sementes consecutivas
        \[
        \Delta\theta=2\pi\Bigl(1-\frac1\varphi\Bigr)=\frac{2\pi}{\varphi^2}\approx 137{,}507^\circ;
        \]
  \item malha erguida uma dimensão --- cone nulo (a espiral plana é a sombra).
\end{enumerate}

As espirais de Fibonacci aparecem sozinhas. Foram precisos: um ponto, o acto de dividir-se,
e a memória de ter dividido.

\subsection{O operador $T$ da recorrência}

O mesmo passo estrutural produz módulo e fase:

\begin{definicao}[operador da espiral de ouro]
\label{def:T-espiral}
\[
T=\varphi\,e^{J\Delta\theta},\qquad
\Delta\theta=\frac{2\pi}{\varphi^2},\qquad
J^2=-1.
\]
\end{definicao}

\begin{teorema}[a espiral resolve a recorrência de onda]
\label{thm:espiral-psi}
Seja $\psi_0=A_0\,e^{J\theta_0}$ com $A_0>0$. A órbita
\[
\psi_{n+1}=T\psi_n
\]
satisfaz
\begin{enumerate}[leftmargin=1.2em,itemsep=1pt]
  \item $|T\psi_n|=\varphi|\psi_n|$ \quad (logo $A_{n+1}/A_n=\varphi$);
  \item $\arg(T\psi_n)-\arg(\psi_n)=\Delta\theta=2\pi/\varphi^2$;
  \item a solução explícita
        \[
        \boxed{\psi_n=A_0\,\varphi^n\,e^{J(\theta_0+n\Delta\theta)}}.
        \]
\end{enumerate}
Não é ansatz: é a órbita de $T$. A recorrência selecciona $A$; o ciclo (ângulo da condição
de máxima simetria) selecciona $\theta$. Este teorema é de consistência da órbita e não
substitui a derivação causal de $T$ a partir de $L_0+L_7$.
\end{teorema}

\begin{proof}
$T$ é produto de escalar positivo $\varphi$ por rotação $e^{J\Delta\theta}$. Logo
$|T\psi|=\varphi|\psi|$ e o argumento avança exactamente $\Delta\theta$. Por indução,
$\psi_n=T^n\psi_0=A_0\varphi^n e^{J(\theta_0+n\Delta\theta)}$.
A identificação $\varphi$ e $\Delta\theta=2\pi/\varphi^2$ é a do catálogo 0.20.2
(máxima simetria sequencial $\Rightarrow\varphi$; giro $\Rightarrow 2\pi/\varphi^2$).
\end{proof}

\subsection{Resíduos}

\begin{center}
\renewcommand{\arraystretch}{1.2}
\begin{tabular}{@{}ll@{}}
$\operatorname{Res}_A$ & $A_{n+1}/A_n-\varphi=0$ \\
$\operatorname{Res}_\theta$ & passo angular $=2\pi/\varphi^2$; fecho plano não periódico
($\varphi$ irracional); fecho no cone nulo (dual) \\
$\operatorname{Res}_\psi$ & os dois simultaneamente na órbita de $T$ \\
\end{tabular}
\end{center}

O resíduo local da espiral,
\[
\operatorname{Res}^{\mathrm{passo}}_\theta
=\Delta\theta-\frac{2\pi}{\varphi^2},
\]
vale zero por construção de $\Delta\theta$. Isso garante a consistência algébrica da etapa angular,
mas não substitui a derivação causal de $\Delta\theta$ a partir da estrutura anterior
($L_0+L_7$ e o catálogo). A dívida é precisamente a passagem de um passo angular bem definido
para a origem causal desse passo.

Há também uma diferença precisa entre o \emph{invariante do par} e o \emph{resíduo de truncagem}.
O par dual mantém
\[
|\psi^+|\,|\psi^-|=1
\]
como conservação do retorno; a truncagem num suporte finito introduz um resíduo de janela
\[
\mathcal R_N=\mathcal T_N(\psi)-\psi,
\]
que mede a perda da cauda e não a quebra do invariante. O que se exige aqui é separar estas duas
coisas: o invariante do par é estrutural, enquanto o erro de suporte é finito e local.
No plano a órbita não fecha (número mais irracional). Uma dimensão acima, o mesmo objecto
é o cone nulo --- onde a distância própria é nula e o crescimento se conserva. O dual da
espiral áurea é o cone nulo (\texttt{tests/cone\_nulo.c}).

\subsection{O que isto fecha e o que deixa para a estatística}

O que fica estabelecido é a forma dinâmica da órbita:
\[
\text{recorrência}\longrightarrow A,\qquad
\text{ciclo}\longrightarrow\theta,\qquad
\Longrightarrow\psi=A\,e^{J\theta}
\]
como solução de $\psi_{n+1}=T\psi_n$, com $T$ lido do catálogo. Isso produz uma
\emph{realização de onda candidata} --- ou seja, a espiral realiza a forma de onda no
sentido dinâmico.

A dívida causal que permanece é a passagem de $L_0+L_7$ para o próprio operador
$T=\varphi e^{J\Delta\theta}$. O teorema prova a consistência da órbita e a forma da solução,
mas não fecha ainda a cadeia causal que produz $T$ a partir das leis anteriores.

Não fecha ainda a leitura estatística $P_U$ dessa $\psi$ (frequências, Born como observação).
A probabilidade observa a realização; não a fabrica. A ordem correta é: realização geométrica,
$\pi_U$, multiplicidade, e só depois $P_U$. O passo que fica para depois do fecho geométrico é a ponte
\[
P_U(x)\stackrel{?}{=}|\psi_x|^2,
\]
que não pode ser usada como hipótese de construção.


% ============================================================
\section{Normalização: transformada, metais e Pisot}
\label{sec:norm}
% ============================================================

A órbita da espiral
\[
\psi_n=A_0\varphi^n e^{J(\theta_0+n\Delta\theta)}
\]
cresce. A folha conjugada (Pisot) contrai:
\[
\bar\varphi=-\frac1\varphi,\qquad
\bigl|\bar\varphi\bigr|^n=\varphi^{-n}\to 0.
\]
Duas folhas da mesma borda $b^2-b-1=0$:
\[
\varphi^n\quad\text{(expansiva)},\qquad
(-\varphi^{-1})^n\quad\text{(contrativa)}.
\]
No limite a contrativa desaparece; a realização física vive na expansiva.
No suporte finito a contrativa ainda contribui residualmente --- e o resíduo tem de
ser controlado.

\subsection{Corte finito e factor da transformada}

Seja $I=\{0,\ldots,N-1\}$ o suporte finito e
\[
\Psi_I=(\psi_0,\ldots,\psi_{N-1})
\]
a órbita truncada. A normalização não é arbitrária: vem da mesma máquina que já
normaliza a transformada (Teor.~0.60; mudança de base):
\[
c=2^{-m}F(f),\qquad N=2^m.
\]
Dois factores, duas leituras:
\begin{center}
\renewcommand{\arraystretch}{1.2}
\begin{tabular}{@{}ll@{}}
$2^{-m}=1/N$ & factor de amplitude da transformação \\
$\sqrt N$ & factor da norma $L^2$ (Parseval: $\|F(f)\|_2^2=N\|f\|_2^2$) \\
\end{tabular}
\end{center}
Os documentos distinguem $N$ de $\sqrt N$ precisamente porque são essas duas leituras,
não dois factores concorrentes. Nos metais / dourada: em $m=0$ o factor reparte-se
($\sqrt N$); em $m=1$ (estrela) junta-se ($N$).

No contínuo / círculo unitário a área traz $\pi$; no discreto de largura $m$ o factor
é potência de 2. As duas escalas não se misturam: cada suporte carrega a sua.

\subsection{Uma dobra, duas leituras}

\begin{center}
Lei 0 + Lei 7 $\;\Longrightarrow\;$ dobra com retorno
\end{center}
A mesma dobra lê-se de dois lados:
\begin{center}
\renewcommand{\arraystretch}{1.25}
\begin{tabular}{@{}ll@{}}
dinâmica & $|\psi^+|\,|\psi^-|=1$ \\
estatística & $G(x)=|\pi_U^{-1}(x)|$ \\
\end{tabular}
\end{center}
$P_U=G/|I|$ não é segunda lei sobre a dinâmica: é a contagem das fibras da mesma projeção.

Ordem:
\[
\text{Lei 0}\to\text{Lei 7}\to\text{vinco}\to\text{retorno}
\to(\psi^+,\psi^-)\to\pi_U\to G\to P_U.
\]

\subsection{Invariante unitário por fibra}

Seja $x$ fixo e
\[
\pi_U^{-1}(x)=\{u_1,\ldots,u_{G(x)}\}
\]
as suas pré-imagens. Cada $u$ carrega o par $(\psi_u^+,\psi_u^-)$. Se o invariante do retorno
for \emph{unitário por pré-imagem},
\[
\boxed{\mathcal I(\psi_u^+,\psi_u^-)=1},
\]
então automaticamente
\[
\sum_{u\in\pi_U^{-1}(x)}\mathcal I(\psi_u^+,\psi_u^-)
=\sum_{u\in\pi_U^{-1}(x)}1
=|\pi_U^{-1}(x)|
=G(x).
\]
A normalização é a contagem:
\[
\boxed{P_U(x)=\frac1{|I|}\sum_{u\in\pi_U^{-1}(x)}1=\frac{G(x)}{|I|}.}
\]

O par não deve ser entendido como duas cópias da mesma estrutura algébrica.
A hipótese de trabalho é mais forte e mais específica: uma face pode realizar
a composição anelar, enquanto a face dual realiza a projeção idempotente.

$$
\boxed{
\psi_u^+\in R^+,\qquad
\psi_u^-\in E^-,
\qquad
(\psi_u^-)^2=\psi_u^-.
}
$$

No lado \(R^+\), a associatividade e a distributividade são

$$
(ab)c=a(bc),
$$

$$
a(b+c)=ab+ac,\qquad
(a+b)c=ac+bc.
$$

No lado \(E^-\), a condição fundamental é

$$
(\psi_u^-)^2=\psi_u^-,
$$

e, caso exista uma composição binária \(\circ\),

$$
(e_1\circ e_2)\circ e_3
=
e_1\circ(e_2\circ e_3).
$$

A dualidade relaciona as duas estruturas por leis de compatibilidade,
não pela identificação dos produtos.

Em particular,

$$
|\psi_u^+|\,|\psi_u^-|=1
$$

é uma conservação de módulo e não deve ser promovida a

$$
\psi_u^+\psi_u^-=1.
$$

Se a segunda igualdade fosse imposta e \(\psi_u^-\) fosse idempotente
e invertível, então necessariamente \(\psi_u^-=1\). A estrutura dual
desapareceria. O papel da dualidade é precisamente evitar essa
identificação: uma face compõe, a outra estabiliza.

Assim, a forma correcta do par é

$$
\boxed{
\text{composição anelar}
\quad\rightleftarrows\quad
\text{projeção idempotente},
\qquad
|\psi^+|\,|\psi^-|=1.
}
$$

O problema estatístico permanece inalterado:

$$
(\psi_u^+,\psi_u^-)
\longrightarrow
\pi_U^{-1}(x)
\longrightarrow
\sum 1
\longrightarrow
G(x)
\longrightarrow
P_U(x).
$$

A nova questão é algébrica: demonstrar que a correspondência entre as
duas faces satisfaz as condições de associatividade internas e as leis
de distributividade/compatibilidade dual, com resíduo zero.

\begin{problema}[invariante unitário por fibra]
\label{prob:norm}
Localizar o invariante $\mathcal I$ do retorno (Lei 0+Lei 7, $|\det|=1$, emparelhamento dual)
tal que
\[
\mathcal I(\psi_u^+,\psi_u^-)=1
\quad\text{para cada }u\in I,
\]
e portanto
\[
\sum_{u\in\pi_U^{-1}(x)}\mathcal I=G(x).
\]
Candidatos já medidos: $|\psi^+||\psi^-|=1$; $\psi^+\psi^-=(-1)^{\mathrm{ord}}$; $|\det|=1$.
Não inventar $\mathcal I$ fora do retorno. Não usar $C_N$, Fourier nem Pisot como fundamento.
\end{problema}

\section{A álgebra dual da realização: anel, idempotente e distributividade}
\label{sec:algebra-dual}

A conservação

$$
\boxed{|\psi_u^+|\,|\psi_u^-|=1}
$$

não exige que as duas faces sejam o mesmo tipo de objecto algébrico.
Ao contrário: a realização pode ser dual precisamente porque as duas faces
carregam estruturas diferentes que se conservam sob a troca.

A distinção fundamental é

$$
\boxed{
\psi^+\ \text{é realizado no lado anelar},
\qquad
\psi^-\ \text{é realizado no lado idempotente}.
}
$$

A dualidade não identifica os dois produtos. Ela relaciona as duas estruturas.

\subsection{Duas faces, duas estruturas}

Seja \(A\) o suporte comum das duas realizações. A face \(+\) é
munida da operação \(\star_+\), enquanto a face \(-\) é munida da
operação transportada \(\star_-\). Seja

$$
\delta:A\to A,\qquad \delta^2=\mathrm{id},
$$

a involução que troca as duas faces. Define-se

$$
\boxed{
 u\star_- v
=
\delta\!\left(\delta(u)\star_+\delta(v)\right).
}
$$

Assim, a segunda face não é uma nova álgebra introduzida
independentemente: é a realização da primeira lei após a troca de
face.

Da definição segue imediatamente

$$
\boxed{
\delta(u\star_+v)
=
\delta(u)\star_-\delta(v).
}
$$

Portanto, em geral,

$$
\delta(u\star_+v)
\neq
\delta(u)\star_+\delta(v),
$$

e isso não representa uma falha de preservação. A operação que deve ser
preservada é a operação da face de chegada.

Em outras palavras,

$$
\boxed{
\delta:(A,\star_+)\xrightarrow{\cong}(A,\star_-).
}
$$

A involução é reversível porque \(\delta^2=\mathrm{id}\). Se
\(\star_+\) é associativa e distributiva, então \(\star_-\) possui as
mesmas propriedades por transporte da estrutura. Por exemplo,

$$
(u\star_-v)\star_-w
=
 u\star_-(v\star_-w),
$$

e as leis distributivas de \(\star_-\) seguem das correspondentes de
\(\star_+\).

O ponto estrutural é, portanto,

$$
\boxed{
\text{mesmo suporte}
+
\text{operação transportada}
=
\text{segunda realização}.
}
$$

A dualidade não exige que \(\delta\) seja automorfismo de uma única
operação. Ela exige que a operação seja trocada juntamente com a face.

\subsection{A dualidade não é inversão algébrica}

O invariante da realização é

$$
\boxed{
|\psi_u^+|\,|\psi_u^-|=1.
}
$$

Esta igualdade é uma igualdade de módulos. Ela não implica

$$
\psi_u^+\psi_u^-=1.
$$

Esta separação é necessária. Se um elemento \(e\) fosse simultaneamente
idempotente e invertível em um anel unitário, então

$$
e^2=e
$$

e, multiplicando pela inversa,

$$
e=1.
$$

Logo,

$$
\boxed{
e^2=e\ \land\ e\text{ invertível}
\Longrightarrow e=1.
}
$$

A face idempotente não deve, portanto, ser interpretada como o inverso
algébrico da face anelar. O exemplo transportado acima mostra precisamente a diferença entre
inversão algébrica e dualidade estrutural. A involução \(\delta\) não
produz, em geral,

$$
\delta(u\star_+v)=\delta(u)\star_+\delta(v).
$$

Ela produz

$$
\boxed{
\delta(u\star_+v)=\delta(u)\star_-\delta(v),
}
$$

isto é, preserva a composição somente depois da troca da operação.

Assim,

$$
\boxed{
\text{dualidade}\neq\text{automorfismo},
\qquad
\text{dualidade}=\text{troca reversível da realização}.
}
$$

Esta distinção é a realização concreta do duomorfismo: a mesma
transformação que não preserva a lei \(\star_+\) como lei fixa é um
isomorfismo entre a estrutura de chegada e a estrutura transportada.

\subsection{Condições mínimas de associatividade}

Para que a composição interna de cada face seja bem definida, as operações
devem satisfazer associatividade no seu próprio sítio.

No lado anelar:

$$
\boxed{
(ab)c=a(bc).
}
$$

No lado idempotente, se \(\circ\) for uma operação binária interna,
exige-se

$$
\boxed{
(e_1\circ e_2)\circ e_3
=
e_1\circ(e_2\circ e_3).
}
$$

A idempotência é então uma condição adicional:

$$
\boxed{
e\circ e=e.
}
$$

Assim, a estrutura mínima do lado \(-\) é um semigrupo idempotente;
se \(\circ\) for comutativa, obtém-se um semilattice quando também houver
as condições correspondentes de ordem.

A associatividade não vem da dualidade. Ela é uma propriedade interna de
cada operação. A dualidade acrescenta as condições que relacionam as duas
faces.

\subsection{Distributividade interna do lado anelar}

No anel \(R^+\), a distributividade é a usual:

$$
\boxed{
a(b+c)=ab+ac
}
$$

e

$$
\boxed{
(a+b)c=ac+bc.
}
$$

Estas duas leis são consequências da estrutura de anel e não devem ser
reintroduzidas como propriedades da face idempotente.

A face idempotente pode possuir uma estrutura distributiva própria, mas
isso deve ser declarado separadamente.

\subsection{Distributividade dual}

A palavra ``dual'' deve designar aqui uma relação entre as operações das
duas faces, e não uma terceira operação.

Se

$$
D:R^+\longleftrightarrow E^-
$$

é a correspondência dual, as operações são compatíveis quando a passagem
de uma face para a outra preserva a forma da composição apropriada.

Uma condição forte é

$$
\boxed{
D(a(b+c))
=
D(ab)\circ D(ac)
}
$$

e, dualmente,

$$
\boxed{
D((a+b)c)
=
D(ac)\circ D(bc).
}
$$

Estas equações não dizem que \(\circ\) seja o mesmo produto de \(R^+\).
Elas dizem que a distribuição realizada no lado anelar possui uma leitura
correspondente no lado idempotente.

Uma forma mais simétrica é introduzir uma operação de ação

$$
\triangleright:R^+\times E^-\to E^-
$$

e exigir

$$
\boxed{
(a+b)\triangleright e
=
(a\triangleright e)\circ(b\triangleright e)
}
$$

e

$$
\boxed{
(ab)\triangleright e
=
a\triangleright(b\triangleright e).
}
$$

Nesse caso, a primeira equação é a distributividade da soma do anel sobre
a composição da face idempotente, enquanto a segunda exprime
compatibilidade da multiplicação com a composição.

A condição dual correspondente pode ser escrita com uma ação

$$
\triangleleft:E^-\times R^+\to R^+
$$

e

$$
\boxed{
e\triangleleft(a+b)
=
(e\triangleleft a)+(e\triangleleft b).
}
$$

As ações não são obrigatórias para a definição do par dual; elas são a
forma adequada de escrever a interação quando as duas faces precisam
operar uma sobre a outra.

\subsection{O que é necessário para o par dual}

A estrutura mínima pode ser resumida em quatro níveis:

\begin{center}
\renewcommand{\arraystretch}{1.3}
\begin{tabular}{@{}lll@{}}
\toprule
\textbf{nível} & \textbf{condição} & \textbf{significado} \\
\midrule
anel & $(ab)c=a(bc)$ & associatividade da face \(+\) \\
anel & $a(b+c)=ab+ac$, $(a+b)c=ac+bc$ & distributividade interna \\
idempotente & $e\circ e=e$ & projeção / estado fixo \\
dual & compatibilidade $D,\triangleright,\triangleleft$ & ligação entre as faces \\
\bottomrule
\end{tabular}
\end{center}

A conservação

$$
|\psi_u^+|\,|\psi_u^-|=1
$$

é então uma condição adicional de fechamento da realização, não uma lei
de associatividade ou distributividade.

\subsection{O par como realização de uma única dobra}

A estrutura completa pode ser representada por

$$
\boxed{
R^+
\ \underset{D^{-1}}{\overset{D}{\rightleftarrows}}\
E^-,
\qquad
|\psi_u^+|\,|\psi_u^-|=1.
}
$$

A face \(+\) realiza composição anelar:

$$
(a,b,c)\mapsto (ab)c=a(bc).
$$

A face \(-\) realiza estabilização:

$$
e\mapsto e\circ e=e.
$$

A troca entre as faces não cria uma terceira entidade. Ela apenas muda a
leitura da mesma realização:

$$
\boxed{
\text{composição}
\quad\rightleftarrows\quad
\text{projeção}.
}
$$

Esta é precisamente a forma adequada para a dualidade da realização:
uma face carrega a dinâmica de composição; a outra carrega o ponto fixo
da projeção.

\subsection{Consequência para o par \((\psi_u^+,\psi_u^-)\)}

Para cada pré-imagem \(u\in I\), associa-se

$$
u\longmapsto(\psi_u^+,\psi_u^-)
$$

com

$$
\boxed{
|\psi_u^+|\,|\psi_u^-|=1,
}
$$

e com as condições estruturais

$$
\boxed{
\psi_u^+\in R^+,
\qquad
\psi_u^-\in E^-,
\qquad
(\psi_u^-)^2=\psi_u^-.
}
$$

A face idempotente não fornece uma segunda amplitude independente.
Ela fornece a realização estabilizada da mesma fibra.

Assim, para uma fibra

$$
\pi_U^{-1}(x)=\{u_1,\ldots,u_{G(x)}\},
$$

cada pré-imagem carrega um par dual e a contagem continua sendo

$$
G(x)=\sum_{u\in\pi_U^{-1}(x)}1.
$$

A dualidade algébrica, portanto, não altera a definição estatística:

$$
\boxed{
P_U(x)=\frac{G(x)}{|I|}.
}
$$

Ela fornece a estrutura que deve explicar por que cada visita possui peso
unitário.

\subsection{Critério de fechamento}

A dívida agora pode ser formulada de modo mais preciso.

Não basta verificar

$$
|\psi_u^+|\,|\psi_u^-|=1.
$$

É necessário verificar simultaneamente:

$$
\boxed{
\begin{aligned}
&(ab)c=a(bc),\\
&a(b+c)=ab+ac,\\
&(a+b)c=ac+bc,\\
&e\circ e=e,\\
&(e_1\circ e_2)\circ e_3
=
 e_1\circ(e_2\circ e_3),\\
&\delta^2=\mathrm{id},\\
&u\star_-v
=
\delta\!\left(\delta(u)\star_+\delta(v)\right),\\
&\delta(u\star_+v)
=
\delta(u)\star_-\delta(v),\\
&|\psi_u^+|\,|\psi_u^-|=1.
\end{aligned}
}
$$

O experimento com a ação escalar e a involução
\(\delta(a,b)=(-a,b)\) realiza explicitamente as duas últimas condições
estruturais de troca: \(\star_-\) é obtida por transporte de
\(\star_+\), e a involução estabelece a correspondência reversível entre
as duas leis. A associatividade e a distributividade da face transportada
não precisam ser postuladas novamente: seguem das mesmas propriedades
da lei de origem.

O que permanece como dívida não é a existência da segunda lei, mas a sua
identificação com a face idempotente \(E^-\) usada na realização de onda.

A dualidade é a reversibilidade entre as duas faces da mesma realização.
O teste de transporte mostra que essa reversibilidade pode trocar a lei
de composição sem destruir as propriedades estruturais ou a norma:

$$
\boxed{
(A,\star_+)
\underset{\delta}{\overset{\delta}{\rightleftarrows}}
(A,\star_-),
\qquad
\star_-
=
\delta\circ\star_+\circ(\delta,\delta).
}
$$

Portanto,

$$
\boxed{
\text{a dualidade preserva o invariante,
mas pode trocar a lei de composição.}
}
$$

Não se introduz uma terceira álgebra. Há um suporte, duas leituras
operacionais e uma involução reversível que as relaciona.

\subsection{Consequência para Born}

Se o invariante por fibra for unitário, então $|\psi|^2$ (na carta analítica) é \emph{leitura}
desse invariante --- não a regra que define a probabilidade. Born deixa de ser axioma:
é a forma analítica de contar fibras de peso 1.

Cadeia:
\[
\underbrace{\text{retorno}}_{\text{Lei 0+7}}
\to
\underbrace{\text{invariante unitário por fibra}}_{\mathcal I=1}
\to
\underbrace{\text{multiplicidade}}_{G=|\pi_U^{-1}|}
\to
\underbrace{\text{probabilidade}}_{P_U=G/|I|}.
\]

\begin{teorema}[conservação dual --- realização]
\label{thm:par-dual}
Na realização áurea, $|\psi^+_{n+1}|\,|\psi^-_{n+1}|=|\psi^+_n|\,|\psi^-_n|$
($\varphi\cdot\varphi^{-1}=1$). Em $N=4$: $|\psi^+||\psi^-|=1$,
$\psi^+\psi^-=(-1)^n$, $|\det|=1$. Realiza Lei~4 e Lei~7; não as funda.
\end{teorema}

\subsection{O que não se faz}

\begin{itemize}[leftmargin=1.2em,itemsep=1pt]
  \item não começar por $(\psi^+,\psi^-)$ e encaixar leis depois;
  \item não normalizar a folha expansiva por $C_N$ para forçar Born;
  \item não inventar $\mathcal I$ / $\mathcal D$ fora do retorno e de $\pi_U$;
  \item não subir $N$ enquanto $\mathcal I=1$ por fibra não estiver
        localizado com resíduo zero nos documentos.
\end{itemize}


\subsection{Estatística fechada: $\mu=1$ por visita}

A conservação estatística deixou de ser aposta. Teor.~0.50:
\[
\mu=\zeta^{-1}\quad\text{vale 1 em cada visita},\qquad
\sum_{u\in\pi^{-1}(x)}\mu(u)=\sum_{u\in\pi^{-1}(x)}1=G(x),
\]
\[
\boxed{P_U(x)=\frac{G(x)}{|I|}}.
\]
A unidade elementar da contagem \emph{já existe}. Não se fabrica densidade por $C_N$
nem por Born. O problema deixa de ser «como obter uma probabilidade» e passa a ser:
\[
\boxed{\text{como a realização }(\psi^+,\psi^-)\text{ produz as fibras de }\pi_U?}
\]
Uma vez produzida a fibra, a estatística está pronta:
\[
(\psi^+,\psi^-)\;\longrightarrow\;\pi_U^{-1}(x)\;\longrightarrow\;
\underbrace{\sum 1}_{G(x)}\;\longrightarrow\;\frac{G(x)}{|I|}.
\]
A última seta, da estatística $P_U$ para a forma analítica $|\psi|^2$, não é usada na construção:
\[
\boxed{P_U(x)\stackrel{?}{=}|\psi_x|^2.}
\]
Essa igualdade é a ponte de Born; não é a definição de $P_U$.

\subsection{Dois fechamentos independentes}

\begin{center}
\renewcommand{\arraystretch}{1.3}
\begin{tabular}{@{}ccc@{}}
composição & $\longrightarrow$ & $\Delta_N=0$ \\
visita & $\longrightarrow$ & $\mu=1$ \\
\multicolumn{3}{c}{$\Downarrow$} \\
invariante por fibra & $\longrightarrow$ & $G(x)$ \\
\end{tabular}
\end{center}

Não se identifica $\Delta_N$ com $\mu$. Encontram-se na projeção.

\subsection{Cadeia e o que resta}

\[
\boxed{
L_0\to L_7\to H\times H^*\to\text{composição}
\to
\begin{cases}\Delta_N=0\\ \mu=1\end{cases}
\to\pi_U\to G\to P_U
}
\]

\begin{problema}[selecção pelas Leis 0+7]
\label{prob:defeitos}
Demonstrar, a partir apenas de Lei~0 e Lei~7, que a realização selecciona
$\Delta_N=0$ e $\mu=1$ por visita. Ordem de ataque:
\begin{enumerate}[leftmargin=1.2em,itemsep=1pt]
  \item Lei~0: operação nos extremos;
  \item Lei~7: como produz $H\times H^*$ sem fundir;
  \item composição: coeficientes impostos pelo par;
  \item fechamento $\Delta_N=0$ (não postulado);
  \item visita: por que $\mu=\zeta^{-1}$ dá 1 (já Teor.~0.50 na ordem);
  \item realização: como $(\psi^+,\psi^-)$ corresponde a pré-imagens de $\pi_U$.
\end{enumerate}
\end{problema}

\subsection{O que Alonzo já localiza}

Catálogo Def.~0.8 / Física Def.~0.133: corpo da \emph{composição} (iterar $f\circ f\circ\cdots$).
\begin{itemize}[leftmargin=1.2em,itemsep=1pt]
  \item dois golden: cantor $=\varphi$, julia $=\psi$, raízes de $x^2=x+1$;
  \item inversor $x^\dagger=-1/x$; produto $\varphi\cdot\psi=-1$ (estaca);
  \item $H\times H^*$: produto componente a componente; \emph{duas normas}, uma por
        tecido, nada se mistura; associa em $4096/4096$;
  \item três dobras da composição $\to$ octonião dual (grau 8, dois tecidos);
  \item vinco $1=\mathrm{cl}(L_0,L_7)$ é Lei~7; $\pi\circ\pi=\pi$ (idempotência).
\end{itemize}
O par $(\varphi,\psi)$ com produto $-1$ e duas normas separadas é a realização dual
já escrita em Alonzo --- não inventada na Parte XVI.

\subsection{Separação em letras garrafais}

\begin{center}
\bfseries conservação não é $P_U$
\end{center}
A conservação produz o invariante unitário; a projeção produz a multiplicidade;
a normalização produz a probabilidade. Não se toma $|\widehat\psi|^2=G/|I|$ como ponto
de partida.

\textbf{Fechado:} $G=|\pi^{-1}|$ (Def.~0.23); $\mu=1$ por visita (Teor.~0.50);
$P_U=G/|I|$; par Alonzo $(\varphi,\psi)$ com $\varphi\cdot\psi=-1$ e duas normas;
$|\psi^+||\psi^-|=1$ na órbita; $\Delta_N=0$ no caso Hurwitz $p=r=t=s=1$.

\textbf{Aberto (dívida causal):}
$L_0+L_7\stackrel{?}{\Longrightarrow}\{\Delta_N=0,\,\mu=1\}$ como \emph{uma} cadeia;
como $(\psi^+,\psi^-)$ produz as fibras de $\pi_U$; e a ponte final
$P_U(x)\stackrel{?}{=}|\psi_x|^2$.

\section{A medição como variável aleatória sobre a onda}
% ============================================================

\begin{teorema}[a medição como variável aleatória sobre $P_U$]
\label{thm:medicao-va}
Seja $(\Omega,\mathcal F,P_U)$ o espaço de probabilidade associado à realização e seja $X:\Omega\to\mathbb R$
uma variável aleatória sobre esse espaço. Dado um conjunto de outcomes $\Omega=\{x_1,\ldots,x_n\}$,
\[
\mathbb E[X]=\sum_{x\in\Omega} X(x)\,P_U(x).
\]
Se $\{P_i\}$ é uma família de projectores ortogonais sobre o espaço espectral (modo de Partículas,
Def.~0.194), com $\sum_i P_i=\mathrm{Id}$, e $A=\sum_i\lambda_i P_i$, então
\[
\mathbb E[X]=\operatorname{Tr}(\rho A)
\]
quando $X$ toma o valor $\lambda_i$ no outcome $i$ e $\rho$ é o estado já medido em Partículas.
A ponte com $|\psi|^2$ não entra na definição da medida: ela fica como objectivo posterior.
\end{teorema}

\begin{proof}
Normalização e não-negatividade seguem de $P_U\ge 0$ e $\sum_x P_U(x)=1$.
A esperança é a definição de $\mathbb{E}[X]$; o traço expande-se na base espectral.
O ponto novo é a leitura: a família de pesos $P_U$ não é postulado de Born; é a contagem
normalizada da realização, já construída como $G/|I|$.
\end{proof}

\begin{proposicao}[ponte com $|\psi|^2$]
\label{prop:born-ponte}
A igualdade
\[
P_U(x)=|\psi_x|^2
\]
não é usada para construir $P_U$. Ela é a ponte de Born, que deve ser provada depois da
realização geométrica e da identificação do invariante unitário por fibra.
\end{proposicao}

\begin{observacao}[Born como consequência estatística]
\label{obs:born-nao-axioma}
Se a ponte da fase fechar com $\operatorname{Res}_\theta=0$ nos sítios de aplicação,
Born deixa de ser postulado e passa a ser a estatística da realização de onda:
a frequência relativa de outcomes aproxima $|\psi_i|^2$.
Não se funde com Parseval da dourada (Obs.~0.220) nem com GLH (Obs.~0.221).
Parseval conserva $\|F(G)\|_2$; a realização de onda levanta $P_U$ a $\psi$.
Operações distintas; sítios distintos.
\end{observacao}

% ============================================================
\section{Independência, produto e frequência}
% ============================================================

Duas variáveis aleatórias sobre o mesmo $(\Omega,\mathcal{F},P)$ são independentes se as
frequências conjuntas factorizam. O produto de espaços realiza a independência de cópias:
a medida produto é novamente contagem normalizada sobre o rectângulo --- outra vez Hurwitz.

A lei dos grandes números, no suporte finito, afirma que a média empírica $S_n/n$ satisfaz
\[
P\bigl(\bigl|S_n/n-\mathbb{E}[X]\bigr|>\varepsilon\bigr)\to 0
\quad(n\to\infty).
\]
Isto justifica a interpretação frequentista: a frequência relativa de outcomes aproxima
$|\psi_i|^2$ em ensaios repetidos sobre cópias no mesmo estado.

% ============================================================
\section{O que esta parte herda e o que não promove}
% ============================================================

\textbf{Herda:}
\begin{itemize}[leftmargin=1.2em,itemsep=1pt]
  \item a realização e a multiplicidade $G$ (Def.~0.23) --- como registo da dobra, não como fonte do módulo;
  \item a contagem $p=G/|I|$ (Teor.~0.44; Obs.~0.214) --- como estatística observável;
  \item os degraus da fase e $\operatorname{Res}_\theta=0$ (Teor.~0.22) --- lado rotacional;
  \item o rotor $x^2=-1$, a estrela $x^2=x+1$, o idempotente $x^2=x$;
  \item a costura exterior+rotor (norma $a^2\to a^2+b^2$);
  \item a espiral de ouro e o seu dual cone nulo (catálogo 0.20.2);
  \item Alonzo como corpo da composição ($\varphi,\psi$);
  \item Lei 0+7: vinco + retorno ($\mathrm{Res}=0$);
  \item Lei 2 no catálogo: nomeia $f\leftrightarrow f^{-1}$;
  \item dentro das potências, $f'=f^{-1}$ força $\varphi$
        (catálogo §0.64; \texttt{tests/aurea.c});
  \item dualidade de medida Gentil--Hurwitz--Lebesgue (Teor.~cat:0.7).
\end{itemize}

\textbf{Não promove:}
\begin{itemize}[leftmargin=1.2em,itemsep=1pt]
  \item $\sqrt{G/|I|}$ a amplitude geométrica fundamental;
  \item fase e amplitude como fontes independentes coladas a posteriori;
  \item Born a lei de U;
  \item fase $\theta$ importada da mecânica quântica;
  \item identificação de Parseval da dourada com $|\psi|^2$.
\end{itemize}

\vspace{0.8em}
\begin{center}
\textit{estatística fechada ($\mu=1$, $G$, $P_U$). Alonzo: $(\varphi,\psi)$, $\varphi\cdot\psi=-1$, duas normas.\newline
\textbf{conservação não é $P_U$.} Dívida causal: $L_0+L_7\Rightarrow\{\Delta_N=0,\mu=1\}$;
$(\psi^+,\psi^-)\to\pi_U^{-1}(x)$.}
\end{center}

\vspace{0.5em}
\noindent
{\small
prova: Teor.~0.44(3); $d(xy)$; $\exp(tJ)$; catálogo 0.20.2 (espiral);
$T=\varphi e^{J 2\pi/\varphi^2}$; Teor.~\ref{thm:espiral-psi}; §0.64; \texttt{tests/aurea.c},
\texttt{tests/cone\_nulo.c}
\quad|\quad
realiza: $\psi_n=A_0\varphi^n e^{J(\theta_0+n\Delta\theta)}$ como órbita de $T$;
$\operatorname{Res}_A=\operatorname{Res}_\theta=0$ no passo da espiral
\quad|\quad
verifica: ângulo $2\pi/\varphi^2$; dual cone nulo; $P_U$ continua sendo leitura de contagem
\quad|\quad
verifica: \texttt{tests/aurea.c} resíduo 0; Heighway $\neq$ espiral sob o mesmo $|I|$;
norma $a^2\to a^2+b^2$; cone nulo dual da espiral áurea
}

\vspace{1.2em}
\noindent
Esta Parte não reabre $I_0$ de Partículas nem de Quântica.
Não declara Born equivalente a Parseval, a $|\det|=1$, nem a $p=k/N$.
Não promove $\sqrt{G/|I|}$ a amplitude geométrica.
Isola a estatística como leitura observável e aponta a espiral
(projeção + rotor + estrela no mesmo objecto) como candidata interna à realização de onda.
O lado rotacional já tem teste de resíduo nulo. O lado do crescimento (recorrência hiperbólica)
permanece como grau de liberdade que precisa ser fechado com o mesmo critério nos sítios de aplicação
(Alonzo, espiral de ouro, modo espectral), sem confundir a sua realização geométrica com a
ponte estatística $P_U\stackrel{?}{=}|\psi|^2$. Em outras palavras, a construção da órbita e a
contagem por fibra estão fechadas; a identificação final com a regra de Born continua como ponte
aberta a demonstrar, e não como premissa do argumento.

\end{document}






